\documentclass[11pt]{article}
\usepackage[a4paper,margin=1in]{geometry}
\usepackage{amsmath,amssymb,amsthm}
\usepackage{hyperref}

\theoremstyle{plain}
\newtheorem{theorem}{Theorem}[section]
\newtheorem{lemma}[theorem]{Lemma}
\newtheorem{proposition}[theorem]{Proposition}
\newtheorem{corollary}[theorem]{Corollary}
\theoremstyle{definition}
\newtheorem{definition}[theorem]{Definition}
\theoremstyle{remark}
\newtheorem{remark}[theorem]{Remark}

% ------------------------------------------------------------------
% Notation, following the companion notes (Guard's Basic Recursive
% Arithmetic, Shoenfield's notation) with code(.) for Goedel encoding.
% ------------------------------------------------------------------
\newcommand{\code}[1]{\mathit{code}(#1)}
\newcommand{\thmT}{\mathbf{th}}
\newcommand{\run}{\mathbf{run}}
\newcommand{\ev}{\mathbf{ev}}
\newcommand{\Fun}{\mathit{Fun}}
\newcommand{\Term}{\mathit{Term}}
\newcommand{\Con}{\mathrm{Con}}
\newcommand{\Thm}[1]{\mathit{Thm}(#1)}
\newcommand{\Kc}{K}
\newcommand{\Kle}{\preccurlyeq}
\newcommand{\nodes}{\mathit{nodes}}
\newcommand{\Berry}{\mathbf{berry}}

\title{A small library on Kolmogorov complexity\\ in Basic Recursive Arithmetic}
\author{}
\date{\today}

\begin{document}
\maketitle

\begin{abstract}
  This note records, in ordinary mathematical notation, the results about
  Kolmogorov complexity that have been machine-checked in the \texttt{T4}
  development (the same Basic Recursive Arithmetic framework as the
  companion papers). They form a small self-contained library: an
  \emph{upper bound} $\Kc(x)=O(\log x)$, a \emph{counting bound}
  $\#\{x:\Kc(x)\le L\}\le 3^{\,L+1}$, the existence of \emph{incompressible}
  numbers and the unboundedness of $\Kc$, an \emph{abundance} refinement,
  \emph{monotonicity}, and finally the \emph{non-computability} of $\Kc$
  by a Berry-program argument. Every statement is proved with no
  postulates and no holes; only the lower bounds use a single consistency
  hypothesis.
\end{abstract}

\section{Setup}

We work in the object theory $T$ of Basic Recursive Arithmetic, with
terms built from $\mathbf{O}$, the unary $\mathbf{s}$ and the
primitive-recursive functionals; the numeral of $n\in\mathbb N$ is
$\overline n=\mathbf{s}^n\mathbf{O}$, and $\Thm{A}$ is the \emph{meta}
judgement ``$A$ is derivable in $T$'' --- in the formalisation, that the
Agda type $\mathrm{Deriv}(A)$ is inhabited. Kolmogorov
complexity~\cite{Kolmogorov65,LiVitanyi} is defined relative to the
universal machine of the companion Chaitin development~\cite{arGII}.

\paragraph{Programs as ternary numbers.}
A \emph{program} is a natural number $p$, read in base $3$; its digits
form a string fed to a fixed parser, so the enumeration of programs is the
identity and ``length $\le L$'' is the bounded condition
\[
   |p|\le L \quad\Longleftrightarrow\quad p < 3^{\,L+1}.
\]
The interpreter is the object function $\run\in\Fun_2$: $\run(\overline
p,\overline N)$ runs program $p$ for $N$ steps. ``Program $p$, with fuel
$N$, describes $x$'' is the provable $\Delta_0$ equation
\[
   \run(\overline p,\overline N)=\mathbf{s}\,\overline x,
\]
the leading $\mathbf{s}$ separating ``halted with value $x$'' from ``not
yet halted''.

\paragraph{The complexity predicate.}
For $L,x\in\mathbb N$ we write $\Kc(x)\le L$ (the relation $\Kle$) for
\[
  \Kc(x)\le L \;\;:\equiv\;\;
  \exists\,p<3^{\,L+1}\;\exists\,N\;\;
  \Thm{\,\run(\overline p,\overline N)=\mathbf{s}\,\overline x\,}\,;
\]
``some program of length $\le L$ provably outputs $x$''. Thus $\Kc(x)$ is
the least description length of $x$. Note that $\Kc(x)\le L$ is an
\emph{existential} over the unbounded fuel $N$, hence not a decidable
predicate; this is why several statements below appear in negative form.

\begin{remark}[A notational caveat: meta versus internal]
The reader should keep in mind that, unlike in the companion
Chaitin--G\"odel~I development, the predicate $\Kc(x)\le L$ here is a
\emph{meta} (Agda) statement --- a meta existential $\exists\,p,N$ over the
\emph{object-derivability} facts $\Thm{\run(\overline p,\overline
N)=\mathbf{s}\,\overline x}$ --- so its negation $\Kc(x)>L$ is likewise a
meta statement, whereas in Chaitin--G\"odel~I $\Kc(x)>L$ is the
\emph{internal} open $\Pi_1$ BRA \emph{formula} $\mathbf{CK}(x,\cdot)\neq
\mathbf{O}$ that $T$ may or may not prove; the two notions agree only
through $\Sigma_1$-completeness, and the shared notation should not be read
as identifying them.
\end{remark}

\paragraph{Reference semantics and consistency.}
Each combinator $f\in\Fun_1$ has a meta value function
$\ev(f):\mathbb N\to\mathbb N$ (the standard primitive-recursive
semantics, against which the universal machine is proved correct). The
only hypothesis used by the lower bounds is simple consistency, taken here
in its \emph{meta} form
\[
  \mathrm{con}\;\equiv\;\neg\,\Thm{\,\mathbf{O}=\mathbf{s}\,\mathbf{O}\,},
\]
the Agda statement that the type
$\mathrm{Deriv}(\mathbf{O}=\mathbf{s}\,\mathbf{O})$ is empty (``$T$ has no
derivation of $0=1$''). This is again a meta statement, \emph{not} the
internal open BRA formula
$\Con_T\equiv(\thmT(\mathbf{x}_0)\neq\code{\mathbf{O}=\mathbf{s}\,\mathbf{O}})$
of the G\"odel~II development --- an object sentence that $T$ itself may or
may not prove.

\section{The upper bound}

\begin{theorem}[Logarithmic upper bound]\label{thm:upper}
For every $x\ge 1$ there are a program $p$ and a fuel $N$ with
\[
  p<3^{\,c+d\cdot D},\qquad
  \Thm{\,\run(\overline p,\overline N)=\mathbf{s}\,\overline x\,},
  \qquad\text{and}\qquad 3^{\,D-1}\le x,
\]
where $D$ is the number of base-$3$ digits of $x$ and $c,d$ are fixed
constants. Consequently $\Kc(x)\le c+d\cdot D = O(\log_3 x)$.
\end{theorem}

The program is the base-$3$ Horner code of $x$: an honest object program
whose size $\nodes$ is linear in $D$, run by the universal-machine
correctness theorem; the bound $p<3^{\,\nodes+1}$ is the tight base-$3$
combinatorial estimate $\,2\cdot p+3\le 3^{\,\nodes+1}$. This is the only
result that needs no consistency hypothesis.

\section{The counting bound}

A program describes at most one value --- on pain of inconsistency, since
two distinct outputs of one program at a common (monotonically aligned)
fuel give $\Thm{\,\mathbf{s}\,\overline{x_i}=\mathbf{s}\,\overline{x_j}\,}$
with $x_i\neq x_j$, hence $\Thm{\mathbf{O}=\mathbf{s}\,\mathbf{O}}$. Counting
programs then bounds the values.

\begin{theorem}[Counting bound]\label{thm:count}
Assume $\mathrm{con}$. There are at most $3^{\,L+1}$ numbers of complexity
$\le L$:
\[
  \#\{\,x:\Kc(x)\le L\,\}\;\le\;3^{\,L+1}.
\]
Constructively: any sequence $x_0,x_1,\dots$ with $\Kc(x_i)\le L$ for all
$i$ has a repeated value among its first $3^{\,L+1}+1$ entries.
\end{theorem}

The constructive form is the honest one: each value picks a describing
program in the interval $\{p<3^{\,L+1}\}$, and the meta pigeonhole forces
two of the first $3^{\,L+1}+1$ indices to share a program, hence (by
determinism) to carry equal values.

\section{Incompressibility and unboundedness}

\begin{theorem}[Incompressible numbers exist]\label{thm:incompress}
Assume $\mathrm{con}$. For every $L$ it is not the case that every $i\le
3^{\,L+1}$ has $\Kc(i)\le L$. Equivalently, some $i\le 3^{\,L+1}$ satisfies
$\Kc(i)>L$.
\end{theorem}

\begin{corollary}[$\Kc$ is unbounded]\label{cor:unbounded}
Assume $\mathrm{con}$. For every $L$ it is not the case that $\Kc(i)\le L$ for
all $i$.
\end{corollary}

There are $3^{\,L+1}+1$ numbers in $[0,3^{\,L+1}]$ but at most $3^{\,L+1}$
of complexity $\le L$, so one of them is incompressible above $L$. Because
$\Kc(x)\le L$ is not decidable, the statement is given in the negative
form ``not all are compressible'' rather than as a pinned witness.

\begin{theorem}[Abundance]\label{thm:abundance}
Assume $\mathrm{con}$. No injective block of $3^{\,L+1}+1$ distinct numbers can
all have complexity $\le L$.
\end{theorem}

Taking the block to be the identity recovers
Theorem~\ref{thm:incompress}; the general form shows incompressible
numbers are unavoidable inside any large collection of distinct values.

\begin{proposition}[Monotonicity]\label{prop:mono}
$\Kc(x)\le a$ and $a\le b$ imply $\Kc(x)\le b$.
\end{proposition}

\section{Non-computability of $\Kc$}

The library culminates in the classical observation that $\Kc$ itself is
not computable~\cite{Chaitin74,LiVitanyi}, by Berry's
paradox~\cite{Russell08}. Say that a total combinator $K\!f\in\Fun_1$
\emph{computes} $\Kc$ if it decides the complexity predicate:
\[
  \ev(K\!f)(x)\le L \quad\Longleftrightarrow\quad \Kc(x)\le L
  \qquad(\text{for all }L,x).
\]
The two directions are soundness ($\Kc(x)\le L\Rightarrow\ev(K\!f)(x)\le
L$, no over-estimate) and realisability ($\ev(K\!f)(x)\le L\Rightarrow
\Kc(x)\le L$, the value is genuinely achieved).

\begin{theorem}[Non-computability]\label{thm:noncomp}
Assume $\mathrm{con}$. No $K\!f\in\Fun_1$ computes $\Kc$.
\end{theorem}

The proof is the Berry paradox, realised as an honest object program.
Fix a hypothetical $K\!f$. Build the fixed, $\mu$-free combinator
\[
  \Berry_{K\!f}\in\Fun_1,\qquad
  \Berry_{K\!f}(L) \;=\; \text{the least } x<3^{\,L+1}+1
                          \text{ with } \ev(K\!f)(x)>L,
\]
a bounded search ``scan the days $0,1,\dots,3^{\,L+1}$ and print the first
incompressible one''. Three facts close the argument, for a suitable $L$.

\begin{itemize}
\item \emph{The search succeeds:} writing $b_L=\Berry_{K\!f}(L)$, one has
  $\ev(K\!f)(b_L)>L$. Indeed if \emph{no} $x\le 3^{\,L+1}$ had
  $\ev(K\!f)(x)>L$, then by realisability every such $x$ would have
  $\Kc(x)\le L$, contradicting Theorem~\ref{thm:incompress}.
\item \emph{$b_L$ has a short description:} the program
  $\Berry_{K\!f}\circ\langle\text{constant }L\rangle$ runs to $b_L$, and
  its size is $\nodes = \text{const}+O(\log L)$ (the constant absorbs the
  fixed $\Berry_{K\!f}$, the logarithm is the base-$3$ Horner code of the
  input $L$). By Theorem~\ref{thm:upper} this gives $\Kc(b_L)\le \nodes$.
\item \emph{Exp beats linear:} since $\nodes=\text{const}+O(\log L)$,
  there is an $L$ with $\nodes\le L$; take $L=3^{\,k}$ for $k$ large, where
  the base-$3$ length of $L$ is $k+1$, so a constant plus a multiple of
  $k+1$ stays below $3^{\,k}$.
\end{itemize}
For that $L$ we have simultaneously $\Kc(b_L)\le \nodes\le L$ (describable
within length $L$) and $\ev(K\!f)(b_L)>L$. Soundness turns the first into
$\ev(K\!f)(b_L)\le L$, contradicting the second. \qed

\begin{remark}
The whole development is checked in Agda under
\texttt{-{}-safe -{}-without-K -{}-exact-split}, with no postulates and no
holes. Theorem~\ref{thm:upper} is consistency-free; every other result
uses $\mathrm{con}$ as its sole hypothesis, through the single fact that one
program describes one value. The reference semantics $\ev$ is connected to
the object machine by a small \emph{derived lemma} --- not a new axiom or a
reflection scheme: from an object equation
$f(\overline k)=\overline{\varphi(k)}$, the already-proved soundness of the
machine, and numeral injectivity ($\overline m=\overline n\Rightarrow m=n$,
itself a consequence of $\mathrm{con}$), one \emph{concludes} the meta equation
$\ev(f)(k)=\varphi(k)$. This is what lets the Berry program be reasoned
about value-by-value without ever unfolding the Cantor pairing of the
universal machine.
\end{remark}

\begin{thebibliography}{10}

\bibitem{Chaitin74}
G.~J. Chaitin.
\newblock Information-theoretic limitations of formal systems.
\newblock {\em Journal of the ACM}, 21(3):403--424, 1974.

\bibitem{arGII}
T.~Coquand.
\newblock Autoformalisation of G\"odel's second incompleteness theorem in
  Basic Recursive Arithmetic.
\newblock arXiv:2606.01898, 2026.
\newblock \url{https://arxiv.org/abs/2606.01898}.

\bibitem{Kolmogorov65}
A.~N. Kolmogorov.
\newblock Three approaches to the quantitative definition of information.
\newblock {\em Problems of Information Transmission}, 1(1):1--7, 1965.

\bibitem{LiVitanyi}
M.~Li and P.~Vit\'anyi.
\newblock {\em An Introduction to Kolmogorov Complexity and Its
  Applications}.
\newblock Springer, 3rd edition, 2008.

\bibitem{Russell08}
B.~Russell.
\newblock Mathematical logic as based on the theory of types.
\newblock {\em American Journal of Mathematics}, 30(3):222--262, 1908.
\newblock (Berry's paradox, attributed to G.~G.~Berry.)

\end{thebibliography}

\end{document}
